\chapter{Materiais e Métodos}\label{cap:materialemetodos}

Este capítulo descreve os dados, a reconstrução dos escores de pipelines e a avaliação dos metamodelos. O protocolo distingue os dados brutos usados na validação dos classificadores da meta-descrição da tarefa rotulada. A reconstrução integral substitui a produção histórica de escores na avaliação principal; os artefatos anteriores são preservados para comparação controlada.

\section{Materiais e Fonte dos Dados}\label{sec:materiais}

Os dados provêm de 201 tarefas OpenML \cite{Vanschoren2014} do acervo original. Os metadados oficiais identificam tarefa, versão do dataset e alvo; os arquivos brutos ficam no cache local OpenML, e seus hashes integram o catálogo. Utiliza-se o alvo especificado pela tarefa, confirmado como igual ao padrão do dataset nesta coleta. Não houve falha de carregamento.

Os critérios de inclusão são classificação supervisionada, 500 a 10.000 exemplos, ao menos duas classes, pelo menos cinco observações por classe e atributos tabulares disponíveis após remover colunas de data/hora. Os critérios não usam desempenho dos classificadores ou do ranqueamento. Permanecem 193 tarefas elegíveis em 165 grupos; o catálogo registra as oito exclusões por suporte de classe insuficiente, nomes, links, tamanho, atributos, classes, balanceamento e proporção de valores ausentes.

Tarefas distintas podem representar o mesmo dataset ou versões relacionadas. Os grupos anteriores, formados por nomes repetidos ou atributos processados idênticos, são preservados e conferidos contra IDs e fingerprints dos atributos brutos. A atribuição de grupos não mudou nesta reconstrução. Essa regra é conservadora, mas não assegura detectar todas as famílias com nomes e transformações diferentes. A tarefa 363713 já integrava o acervo e não é tratada como holdout externo.

\subsection{Ambiente Computacional}\label{subsec:recursos}

A execução usa Python 3.14.2, NumPy 2.4.3, pandas 3.0.1, SciPy 1.17.1 e scikit-learn 1.8.0. O computador local dispõe de processador Intel Core i5-9400F, com seis núcleos lógicos. A reconstrução utiliza quatro processos inicialmente e seis após retomada dos checkpoints em tarefas diferentes, com uma thread BLAS e um job por classificador em cada processo. As versões efetivas, configurações, ajustes operacionais e hashes constam dos manifestos; a limitação de paralelismo controla a utilização de recursos sem subamostrar observações.

\section{Fase 1: Escores em Dados Brutos}

Para cada tarefa, uma validação estratificada de cinco folds usa embaralhamento e semente 42. Cada observação participa do teste uma única vez e do treino nos demais folds. Os índices são gerados uma vez por tarefa e compartilhados por todos os pipelines. Não existe subamostra exclusiva da SVM: os cinco classificadores recebem as mesmas linhas de treino e teste em cada fold.

O tratamento numérico estima a mediana exclusivamente no treino, preservando atributos vazios. Nos atributos categóricos, valores ausentes são convertidos deterministicamente em um marcador; um codificador ordinal é ajustado no treino, atribuindo -1 a categorias desconhecidas no teste. Escalonadores e seletores também são ajustados apenas no treino. O teste recebe somente as transformações ajustadas, sem usar seus rótulos para selecionar atributos ou estimar parâmetros.

As transformações comuns são calculadas uma vez por fold e combinação e reutilizadas pelos classificadores que recebem exatamente a mesma entrada. Essa otimização não compartilha ajuste entre folds e não muda os algoritmos. A auditoria verifica cobertura de linhas, ausência de interseção treino/teste e amostras equivalentes entre classificadores.

\subsection{Arsenal de 60 Pipelines}

O arsenal cruza quatro escalonadores (StandardScaler, MinMaxScaler, RobustScaler e ausência), três seletores (PCA, SelectKBest por ANOVA e ausência) e cinco classificadores. Os seletores usam $q=\min(10,\max(1,P-1),N_{treino})$ componentes ou atributos, respeitando dimensões disponíveis. PCA usa solver automático e semente 42, controlando a parte aleatória não fixada no código histórico.

Os classificadores são: (i) Random Forest com 100 árvores; (ii) SVM RBF com $C=1$, gamma scale e sem limite de subamostra; (iii) k-NN com cinco vizinhos uniformes; (iv) MLP com camadas 64 e 32, otimizador Adam e até 200 iterações; e (v) árvore de decisão com profundidade máxima 10. Usa-se semente 42 nos componentes estocásticos. Demais parâmetros são os padrões da versão registrada da biblioteca, exportados para auditoria. Não se otimizam os hiperparâmetros desse arsenal nesta comparação.

\subsection{Métricas, Avisos e Resultados Ausentes}

São calculados acurácia balanceada e F1 macro no teste de cada fold. O alvo do metamodelo é a média de acurácia balanceada dos cinco folds, não uma média de subconjunto de folds bem-sucedidos. Se qualquer ajuste falhar, o alvo desse par tarefa/pipeline fica ausente. Médias, desvios amostrais, tempos, estados, erros e avisos são registrados separadamente.

Avisos de convergência, especialmente no MLP, são preservados como evidência do orçamento finito de treinamento. Eles não são ocultados nem descritos como convergência garantida. Os escores resultantes caracterizam os modelos com o orçamento fixado. No total, a execução contempla 57900 ajustes para 11580 posições de escore na matriz final.

\section{Representação Estatística das Tarefas}\label{sec:representacaoReconstruida}

A representação é um vetor determinístico, sem treinamento neural ou otimização de função objetivo. Para cada atributo numérico após codificação da tarefa, calculam-se média, desvio padrão populacional acrescido de $10^{-8}$, mínimo, máximo, mediana, assimetria, curtose de Fisher e correlações de Pearson e Spearman com o alvo codificado. Cada um dos nove descritores é agregado entre atributos por mínimo, percentil 25, média, percentil 75 e máximo: $9\times5=45$ componentes.

Dez propriedades globais completam o vetor: $\log(1+N)$, $\log(1+P)$, número de classes, proporções majoritária e minoritária, entropia das proporções acrescidas de $10^{-8}$, proporção de atributos constantes, média da correlação absoluta entre atributos com diagonal zerada, média das proporções de valores únicos e $P/(N+10^{-8})$. A média da matriz de correlação inclui os zeros da diagonal no denominador $P^2$. Concatenam-se 55 descritores e nove zeros; não há truncamento com dimensão 64. Valores não finitos recebem \texttt{nan\_to\_num} e aplica-se norma L2, preservando vetor nulo.

A meta-descrição utiliza a tarefa rotulada completa, disponível no cenário de recomendação, e é calculada separadamente do ajuste dos classificadores em folds. Ela não usa predições ou escores de teste. Portanto, o estudo não cobre consulta sem rótulos e não trata essa extração global de descritores como uma transformação do classificador dentro da CV. As correlações com alvos nominais multiclasses dependem de codificação arbitrária. O vetor estatístico é distinto de uma representação neural aprendida como Dataset2Vec \cite{Jomaa2021}.

\section{Fase 2: Predição e Avaliação por Grupos}

O k-NN estima a acurácia dos 60 pipelines pela média dos vizinhos ponderada pelo inverso da distância euclidiana:
\begin{equation}\label{eq:knn}
\hat A(P_j)=\frac{\sum_i w_i A_{i,j}}{\sum_i w_i},\qquad
w_i=\frac{1}{\lVert\mathbf{x}_{novo}-\mathbf{x}_i\rVert_2+10^{-8}}.
\end{equation}

O comparador Random Forest meta é um regressor multi-saída nativo com 200 árvores, mínimo de duas observações por folha, max features 1,0 e semente 42. As baselines são média global de escores por pipeline apenas no treino e ranking aleatório por escores uniformes com semente 42 acrescida do fold. Avalia-se também k-NN com padronização de dimensões ajustada exclusivamente no treino.

Usa-se GroupKFold externo com cinco divisões, shuffle e semente 42. Dentro do treino externo, quatro divisões por grupos selecionam $k\in\{1,3,5,7,10\}$ pela média de Kendall ponderado por grupo, escolhendo o menor k em empate. Teste externo não participa da escolha de k, padronização ou imputação de alvos. As divisões têm grupos disjuntos em todos os níveis.

Resultados ausentes no treino meta são imputados pela mediana de cada pipeline. Colunas inteiramente ausentes no treino são excluídas da avaliação correspondente no teste. Escores reais ausentes no teste nunca são imputados para calcular qualidade do ranking. Correlações indefinidas permanecem como NaN com estado explícito, mantendo NDCG e regret quando calculáveis.

\subsection{Métricas de Ranqueamento}

Spearman mede associação de postos. Kendall tau-b trata empates, com
\begin{equation}\label{eq:kendall}
\tau_b=\frac{C-D}{\sqrt{(C+D+T_x)(C+D+T_y)}},
\end{equation}
onde $C,D$ são pares concordantes e discordantes e $T_x,T_y$ contam empates exclusivos. Kendall ponderado padrão da SciPy usa pesos hiperbólicos aditivos $1/(r_i+1)+1/(r_j+1)$ e média de duas ordenações lexicográficas decrescentes dos escores real e predito; sua implementação não fornece p-valor.

NDCG usa ganho linear de acurácia balanceada e desconto logarítmico, conforme a implementação do scikit-learn:
\begin{equation}\label{eq:ndcg}
\operatorname{NDCG}@K=\frac{\operatorname{DCG}@K}{\operatorname{IDCG}@K},\qquad
\operatorname{DCG}@K=\sum_{i=1}^K\frac{rel_i}{\log_2(i+1)}.
\end{equation}
Reportam-se NDCG completo e NDCG@5. Regret@1 é a diferença para a melhor acurácia real disponível ao recomendar o primeiro pipeline; Regret@5 considera o melhor entre os cinco primeiros. Empates recebem tratamento médio em NDCG e ordem estável de colunas em regret. As métricas são complementares: NDCG elevado não equivale a recuperação exata do top-5.

\subsection{Comparação Controlada e Diagnóstico}

Executam-se três condições com as mesmas tarefas, grupos e folds externos: matrizes históricas; escores reconstruídos com vetores históricos; e escores e vetores reconstruídos. As entradas usam precisão float32 em todas as condições, como no armazenamento histórico, evitando mudar a resolução dos empates somente na condição nova. Os valores originais reconstruídos em float64 são preservados. A comparação pareada examina as mesmas tarefas/grupos avaliáveis nas duas condições. Ela distingue parte do efeito da representação, mas não isola causalmente somente a imputação, pois também mudam a codificação por fold, a amostra da SVM e o controle aleatório da PCA.

Calculam-se médias dentro dos grupos e depois entre grupos. Intervalos descritivos de 95\% usam 2.000 reamostragens de grupos, com semente 42; a sobreposição dos treinos limita a interpretação inferencial. O diagnóstico usa cinco associações de Spearman por grupo (tamanho, atributos, proporção majoritária, atributos constantes e distância ao vizinho), com ajuste Benjamini--Hochberg e sem atribuição causal.

\section{Reprodutibilidade e Delimitação}

Arquivos separados preservam parâmetros, versões, hashes de dados e código, predições, métricas, índices e estados de cada ajuste. Checkpoints por tarefa permitem retomar sem misturar execuções de configurações distintas. Dados brutos permanecem no cache local; as evidências e matrizes agregadas são versionadas.

O estudo continua retrospectivo, com consulta rotulada e arsenal fixo. Não houve busca exaustiva de hiperparâmetros ou coleta de teste externa após o desenvolvimento. O piloto anterior de 180 avaliações permanece separado. Essas limitações delimitam as conclusões, mesmo após concluir a reconstrução integral da Fase 1.
